\section{Initial Detection Attempts}
\label{sec:first_attempts}

% PAGE TARGET: ~3 pages.
% Andrew-style failure narrative: two approaches that did not work or did
% not work well, leading to the revised multi-phi method of Section 3.

Two approaches attempted to extract laid lines from a DRP stack: the first one 
exploited the per-pixel directional information directly and built post-processing on top of it; while that
approach turned out to be biased, a second approach was tried by applying a Fourier estimator on a single post-processed grazing-light image.
That second approach worked on clean folios but proved fragile on the
real benchmark. The two attempts are described here in turn, both
because they shaped the design of the final pipeline in
\Cref{sec:multiphi}.

\subsection{Direct DRP per-pixel analysis}
\label{sec:first_attempts:legacy}

The first detector reduced each pixel's DRP to a dominant azimuth. For pixel $(x, y)$, the
reflectance values $I(x, y, \theta, \phi)$ were first averaged over
elevation $\theta$ to give a one-dimensional profile in azimuth. The
mean was subtracted and the first-order Fourier coefficient of the
residual was taken as a complex number $X + i Y$. The magnitude
$\sqrt{X^2 + Y^2}$ measured how strongly that pixel lies toward one
azimuth, and the angle $\mathrm{atan2}(Y, X)$ gave the dominant
direction. The output of the stage was therefore two maps with the
same spatial size as the folio: one of anisotropy strength and one of
azimuth.

\begin{figure}[!htbp]
  \centering
  \includegraphics[width=\linewidth]{anisotropy_azimuth}
  \caption{Per-pixel DRP-derived direction maps on the benchmark folio
  (dataset~4, downsampled by~8, single-folio crop). Left: anisotropy
  magnitude $\|M\|$, the normalised modulus of the first azimuthal
  Fourier coefficient; the watermark structure of the mould is clearly
  visible. Right: azimuth $\arg(M)$ (cyclic colour scale,
  $\pm180^\circ$); the laid-line direction shows as a banded structure
  but is heavily contaminated by per-pixel noise.}
  \label{fig:fa:dirmap}
\end{figure}

A simple direction map of this kind is too noisy to use as the input
to a frequency detector. The inks, stains and surface unevenness of
real manuscripts produce strong local anisotropies that dominate the
response over the much weaker laid-line signal; the laid wires
themselves are not perfectly straight and aligned, so different
regions of the same folio carry slightly different local directions
and a global detector finds no single dominant pattern. The next
stage therefore tiled the page with overlapping $512 \times 512$ pixel
patches and computed, in each patch, the circular mean direction of
the per-pixel azimuth map. Each pixel then received a Gaussian
likelihood of agreeing with its patch's mean direction, and the patch
responses were blended across the page with a Hann window. The output
was a greyscale image in which the laid-line orientation locally
dominated and which could be fed to a frequency detector.

\begin{figure}[!htbp]
  \centering
  \includegraphics[width=0.85\linewidth]{trigmask}
  \caption{Patchwise trigonometric mask computed from the azimuth map
  of \Cref{fig:fa:dirmap}. Per-patch dominant directions are
  collapsed into a Gaussian likelihood map (
  $\sigma_\phi = 12^\circ$, $96 \times 96$ pixel patches with
  $50\%$ overlap), then stitched across the page with a Hann window.
  The laid-line banding is locally enhanced relative to
  ink, fibre noise, and the structural watermark features still
  visible in the underlying anisotropy.}
  \label{fig:fa:trigmask}
\end{figure}

The frequency detector was a patchwise Gabor bank scanning periods
from 4 to 80 pixels. For each patch, the period at which the
filter response was strongest was taken as the local laid-line
period, and a global density estimate was formed as the
score-weighted vote of the patch winners. The pipeline failed
cleanly on the synthetic phantom of \Cref{sec:benchmark:msi:phantom}.
\Cref{fig:fa:legacy_phantom} sweeps the true phantom period from
$12$ to $40$ pixels and runs both the patchwise Gabor and the
final single-image FFT detector of
\Cref{sec:first_attempts:simple} for reference. The FFT detector
tracks the diagonal within two per cent across the sweep; the
patchwise Gabor estimate collapses to the lower boundary of the
candidate-period grid at four pixels for true periods $16$, $24$,
$28$, $32$ and $40$, with a single high-side excursion to
$41$ pixels at true period $20$.

\begin{figure}[!htbp]
  \centering
  \includegraphics[width=\linewidth]{legacy_vs_simple_phantom}
  \caption{Recovered period from the legacy patchwise Gabor detector
  and from the final single-image FFT detector on the synthetic
  Gaussian-comb phantom (SNR $= 15$, wire $\sigma = 2$ px), swept
  across true periods from $12$ to $40$ pixels. Left: recovered vs
  true period; the FFT detector lies on the diagonal while the
  Gabor estimate is dominated by alias responses at the boundary of
  the candidate-period grid. Right: signed period error.}
  \label{fig:fa:legacy_phantom}
\end{figure}

The failure was traced to two compounding biases in the Gabor stage.
The squared (absolute) response of a Gabor filter doubles the
fundamental frequency of any input that crosses zero, so a
sinusoidal laid-line signal at frequency $f_0$ contributes a strong
peak at $2 f_0$ to the patchwise score. The Gabor bank itself is
constant-Q, with kernel sigma scaling proportionally to the period,
so a higher-frequency response is integrated over a smaller spatial
support and scored more favourably than the true peak. The two
biases pull in the same direction across the candidate grid and
drive the per-patch winner to the smallest available period; on the
sweep of \Cref{fig:fa:legacy_phantom} that period is four pixels.

The lesson from this attempt was straightforward. The DRP carries
useful directional information, but reducing it to per-pixel azimuth
maps before doing frequency analysis throws away too much of the
signal-to-noise on the way through. The revised pipeline goes back
to the raw images and uses the DRP only as a way to obtain several
independent grazing-light views of the same folio.

\subsection{Single grazing-light image with radial FFT}
\label{sec:first_attempts:simple}

The second attempt bypassed the direction map and the patchwise
enhancement altogether. From the DRP stack a single grazing-light
image was extracted at the steepest available elevation, and the
one-dimensional power spectrum of the column-mean of that image
(after rotation so that the laid lines ran vertically) was
computed. The dominant peak in the spectrum, restricted to a
plausible search band of eight to eighty pixels, gave the laid-line
period directly. A single Gabor filter at the detected period was
then applied to clean the one-dimensional signal, the phase was
fitted to find the wire positions, and a grid was overlaid on the
image. Wire diameter was estimated from the harmonic amplitudes of
the broadband signal using a Gaussian-comb model.

This pipeline is short and fast, and on a clean synthetic phantom it
recovers the period to better than one per cent. On real folios it
is also accurate when the chosen grazing image happens to be the one
in which the laid lines are most visible. Its weakness is that the
choice of grazing image is essentially arbitrary: depending on the
folio surface curvature, the azimuth of the chosen view, and the
location of ink coverage, one grazing image can give a clean peak
while another gives mostly noise. The simple pipeline therefore
trades the bias of the legacy approach for a single-frame brittleness
that is hard to predict in advance.

\begin{figure}[!htbp]
  \centering
  \includegraphics[width=\linewidth]{grazing_pair}
  \caption{Two grazing-light views of the same folio acquired at the
  same elevation $\theta = 15^\circ$ but at azimuths $\phi$ separated
  by $90^\circ$ (dataset~4, centre crop). Left ($\phi = 0^\circ$,
  source azimuth normal to the laid lines): strong, directional
  banding is clearly visible. Right ($\phi = 90^\circ$, source azimuth
  parallel to the laid lines): the banding is almost suppressed.
  A single-image radial FFT detector run on the right-hand frame
  alone returns mostly noise, which motivates the multi-$\phi$
  aggregation of \Cref{sec:multiphi}.}
  \label{fig:fa:grazing_pair}
\end{figure}

